\begin{table}[t]\caption{Normalised shape minimum $\hat\nu$ in the census: quantiles, the Hermite bound $\gamma_{d-1}$, medians by type of first witness, and the share of fields with $h^+=h$ (equivalently $U_K^+=U_K^2$).}\label{tab:shape}
\begin{tabular}{rrrrrrrrrrr}\toprule
$d$ & $q_{0.01}$ & $q_{0.1}$ & $q_{0.5}$ & $q_{0.9}$ & max & $\gamma_{d-1}$ & U & T & N & $h^+=h$\\\midrule
4 & 0.08 & 0.12 & 0.58 & 0.90 & 1.19 & 1.26 & 0.55 & 0.58 & 0.62 & 30.7\%\\
5 & 0.47 & 0.55 & 0.73 & 0.94 & 1.20 & 1.41 & 0.73 & 0.72 & 0.80 & 45.5\%\\
6 & 0.27 & 0.34 & 0.74 & 0.95 & 1.32 & 1.52 & 0.74 & 0.77 & 0.37 & 40.6\%\\
7 & 0.85 & 0.88 & 0.97 & 1.08 & 1.25 & 1.67 & 0.96 & 0.98 & -- & 48.2\%\\
\bottomrule\end{tabular}\end{table}

Table~\ref{tab:shape} summarises the shapes. Three features stand out.
(a) In the prime degrees $5$ and $7$ the shape minima concentrate in a band around $0.7$--$1.0$, well
below the Hermite bound, while in degrees $4$, $6$, $8$ there is a heavy lower tail: these are the
imprimitive fields, whose minimal vector lies in a subfield of small variance and whose shape is
therefore skewed. (b) In degree $6$, witnesses of type T (small trace, i.e.\ $t(K)$ small
compared with $\nu$ and $\nu_2$: Corollary~\ref{cor:dich}(4)) occur for well-rounded shapes and
witnesses of type N for skewed ones: the medians of $\hat\nu$ are $0.77$ and $0.37$. This is the
budget--cost mechanism seen statistically: a round shape forces $\nu$ to be large compared with the
cheapest totally positive element, while a skewed (imprimitive) shape inherits a small $\nu$ from a
subfield and only a norm-type witness is left. In degree $4$ the distinction is not visible.
(c) Between one half and two thirds of the fields have $U_K^{+}\neq U_K^{2}$; in odd degree all of
these are excluded by the unit criterion, since then $(2)$ is not the square of an ideal and $2u$
can never be a square. The proportion of fields with $h^{+}=h$ that we observe ($31$--$48\%$ in
degrees $4$--$7$) is the kind of statistic predicted by the heuristics of Dummit and Voight for
signature ranks of units \cite{DV}; we have not compared the numbers with their predictions.
